ここでは、個別の高被引用論文が長期的に影響を持つかどうかを検証するため、その寿命構造を分析する。まず、オブソレッセンス状態への到達を判定するために、P.D.B Paroloら\cite{Par1}の長期被引用モデルおよびD. Wangら\cite{Wan1}の被引用減衰モデルを前提とし、次の3つの条件を設定する： 1. 出版年から十分な期間がある、2. ピーク期が十分に過ぎ去っている、3. 被引用が十分に減衰している。これら3つの条件の具体的な設定値を図\ref{fig:gr6}に示す。この判定の結果、536論文中77論文がオブソレッセンス状態にあると判定された。オブソレッセンス状態にあると判定された高被引用論文および高被引用論文全体の、各指標値を図\ref{fig:arr1}に示す。オブソレッセンス状態にあると判定された論文では出版からピークまでの期間のモードは4年、ピークから被引用の十分な減衰までの期間のモードは7年、出版から被引用の十分な減衰までの期間のモードは12年であった。この結果は、P.D.B Paroloらの報告（生物分野では出版後5年以内にピークが訪れる）を支持している。一方、高被引用論文全体では出版からピークまでの期間は10年付近であり異なる様相を示している。オブソレッセンス状態にあると判定された高被引用論文における被引用数のプロットを図\ref{fig:gr7}に示す。出版から被引用が十分に減衰するまでの期間は最小で10年、最大で84年であった。同図からは近年の論文ほど出版直後に高い被引用を受ける傾向があると推察される。
次に、オブソレッセンス状態にあると判定された高被引用論文について、その群内で特徴差があるのかを確認するため階層クラスタリングを行った。クラスタリングに用いた特徴量は、a. 経過年に関する指標としてa-1:出版からピークまでの期間（Peak period）、a-2:ピークから被引用の十分な減衰が見られるまでの期間（Decay period）、b.被引用の寄与に関する指標としてb-1:ピーク年の被引用数（Peak count）、b-2:減衰年までの被引用数合計（Decay cited: A+B）、c.被引用の安定性に関する指標としてc-1:毎年の被引用数の分散（Variance:cited）、c-2:毎年の被引用速度の分散（Variance:diff）、c-3:毎年の引用加速度の分散（Variance:diff-accel）（以上、7指標）である。距離にはユークリッド距離を用い、各指標値にはノーマライズを行った。分散の計算においては先頭に0-パディングを行った5年平均を用いた。この結果を図\ref{fig:gr8}に示す。クラスタリングではノードが集中的に存在する領域といくつかのノードがまばらに見られる周辺領域が観察されるが、明確な分割構造は確認できなかった。
また、表\ref{tab:mostcitedA}に示される最多被引用論文はいずれもオブソレッセンス状態にあると判定されなかった。これら3件の論文における、各指標値を表\ref{tab:mostcitedObs}に、被引用数の経年プロットを図\ref{fig:nonobs}に示す。いずれの論文も出版からピークまでの期間が非常に長く、かつ、いまだに毎年十分な数の被引用がある。うち2件の論文（Protein measurement with ... と Cleavage of structural proteins ...）については多峰性が見られ、被引用の傾向も似通っている。
